\begin{figure*}[t]
\centering
\includegraphics[width=\textwidth]{full/Figure_B1_fixed_task1_gradient_full.pdf}
\caption{Fixed initial Task-1 descent directions and their effect on Task-0 performance. For each method m and training seed s, compute the mean Task-1 cross-entropy gradient at that method/seed's own Task-0 checkpoint over all 997 Task-1 training images, using the exactly reconstructed, untrained new-20-class head. Normalize its negative once and keep this direction fixed throughout the radius scan. The initial new-class head is shared by the two methods within each seed; the resulting directions are method-specific. Every point starts from the same corresponding Task-0 anchor, with all perturbations in the 24 effective Q/V matrices. The original Task-0 classifier is retained for every displayed measurement. Left: signed Task-0 training loss change $D_{0,m,s}(r)=L_{0,m,s}(\theta_{m,s}(r))-L_{0,m,s}(\theta_{m,s}(0))$, averaged over all 1,259 Task-0 training images. The overbar is the within-method arithmetic mean over three seeds. Middle: Task-0 training accuracy (\%), on the same 1,259 images. Right: raw empirical Fisher trace recomputed at each position on the fixed 256-example Task-0 training subset, as the mean squared norm of per-example Task-0 true-label CE gradients. No additional Fisher normalization or multiplication by lambda is applied. Curves and shading show the arithmetic mean and one sample standard deviation (ddof=1) across seeds 0, 1, and 2; a standard-deviation band may extend below zero without implying negative observed Fisher values. The full figure contains all 93 measured radii up to 50, and the companion shows the measured 0–2 interval. Connecting lines are guides between measured points. Both use formal ImageNet-A, ViT-B/16, rank-10 Q/V LoRA. This is a fixed initial CE direction in effective weights, not an actual sequence of optimizer steps or later-task checkpoints. Both the anchor and direction are method-specific, so the comparison does not isolate the anchor alone. Tiny-step descent is verified for Task-1; neither Task-0 damage nor large-radius Task-1 monotonicity is prescribed.
}
\label{fig:task1-direction-full}
\end{figure*}

\begin{figure*}[t]
\centering
\includegraphics[width=\textwidth]{near_origin/Figure_B1_fixed_task1_gradient_near_origin.pdf}
\caption{Fixed initial Task-1 descent directions and their effect on Task-0 performance. For each method m and training seed s, compute the mean Task-1 cross-entropy gradient at that method/seed's own Task-0 checkpoint over all 997 Task-1 training images, using the exactly reconstructed, untrained new-20-class head. Normalize its negative once and keep this direction fixed throughout the radius scan. The initial new-class head is shared by the two methods within each seed; the resulting directions are method-specific. Every point starts from the same corresponding Task-0 anchor, with all perturbations in the 24 effective Q/V matrices. The original Task-0 classifier is retained for every displayed measurement. Left: signed Task-0 training loss change $D_{0,m,s}(r)=L_{0,m,s}(\theta_{m,s}(r))-L_{0,m,s}(\theta_{m,s}(0))$, averaged over all 1,259 Task-0 training images. The overbar is the within-method arithmetic mean over three seeds. Middle: Task-0 training accuracy (\%), on the same 1,259 images. Right: raw empirical Fisher trace recomputed at each position on the fixed 256-example Task-0 training subset, as the mean squared norm of per-example Task-0 true-label CE gradients. No additional Fisher normalization or multiplication by lambda is applied. Curves and shading show the arithmetic mean and one sample standard deviation (ddof=1) across seeds 0, 1, and 2; a standard-deviation band may extend below zero without implying negative observed Fisher values. The full figure contains all 93 measured radii up to 50, and the companion shows the measured 0–2 interval. Connecting lines are guides between measured points. Both use formal ImageNet-A, ViT-B/16, rank-10 Q/V LoRA. This is a fixed initial CE direction in effective weights, not an actual sequence of optimizer steps or later-task checkpoints. Both the anchor and direction are method-specific, so the comparison does not isolate the anchor alone. Tiny-step descent is verified for Task-1; neither Task-0 damage nor large-radius Task-1 monotonicity is prescribed.
}
\label{fig:task1-direction-near-origin}
\end{figure*}
